\usepackage[utf8]{inputenc}
\usepackage[french, arabic, main=english]{babel}
\usepackage{csquotes}
\usepackage[LAE, T1]{fontenc}
\usepackage{array}
\usepackage{amsmath, amssymb, amsfonts, amsthm}
\usepackage{bbm}
% Police principale : Latin Modern (texte + maths), pour matcher exactement le
% rapport de référence (Mezenner & Yekene), qui utilise LMRoman/LMMono/LMMath.
\usepackage{lmodern}
% ---------------------------------------------------------------------------------------
\usepackage[dvipsnames,svgnames]{xcolor}
\usepackage{fancyhdr}
\usepackage{graphicx}
\graphicspath{{assets/images}{assets/pdf}{assets/python}}
\usepackage[inline]{enumitem}

% Espacement de listes par défaut (enumitem) : reproduit la mise en page classique
% de la référence (paragraphes à alinéa, sans parskip).
\setlist[enumerate,1]{label=\arabic*.}

% \paragraph laissé à sa définition par défaut (titre gras en début de ligne,
% « run-in ») : c'est exactement le style des intertitres de la référence.
\usepackage{import}
\usepackage{hyperref}
\hypersetup{%
  colorlinks=true,
  allcolors=black,
  pdfauthor={Djihene Guitoun},
  pdftitle={A Global Framework for GenAI-Enabled Multi-Agent Semantic Communication and Network Adaptation},
  pdfkeywords={GenAI, multi-agent, semantic communication, network adaptation}
}
% cleveref est chargé en toute fin de préambule (voir plus bas) : sinon un paquet
% chargé après lui (glossaries-extra, bookmark, …) re-patche \@chapter/\@part et
% casse la capture du type des labels, ce qui fait rendre \Cref{part:...} en « ?? ».
\usepackage[top=2.5cm,bottom=2.5cm,right=2.5cm,left=2.5cm]{geometry}
\usepackage{appendix}
\usepackage{wrapfig}
\usepackage{cprotect}
\import{preambule}{tikz}

\newlength\figureheight
\newlength\figurewidth

\usepackage[backend=biber, style=numeric-comp, sorting=none]{biblatex}
\addbibresource{references.bib}

\usepackage{etoolbox}

\usepackage{caption}

\import{preambule}{codestyle}
%% Algorithm style
\usepackage[ruled,linesnumbered,nosemicolon]{algorithm2e}
\SetFuncArgSty{}
\definecolor{codegreen}{rgb}{0,0.6,0}
\definecolor{codegray}{rgb}{0.5,0.5,0.5}
\newcommand{\mycomntsty}[1]{\texttt{\textit{\textcolor{codegray}{#1}}}}
\newlength{\commentWidth}
\setlength{\commentWidth}{7cm}
\SetCommentSty{mycomntsty}
\newcommand{\atcp}[1]{\tcp*[r]{\makebox[\commentWidth]{#1\hfill}}}
% no line break version of \atcp
\newcommand{\atcpl}[1]{\tcp*[f]{\makebox[\commentWidth]{#1\hfill}}}

\SetKwInput{KwRequire}{Require}
\SetKwInput{KwEnsure}{Ensure}
\SetKw{Continue}{continue}
% arabtex écrase \Return : utiliser \KwRet dans les algorithmes.

\usepackage{arabtex}
\usepackage{utf8}
\setcode{utf8}


\makeatletter
\def\@fileswithoptions#1{%
\@ifnextchar[%]
{\@fileswith@ptions#1}%
{\@fileswith@ptions#1[]}}
\makeatother


\usepackage{afterpage}
\usepackage{float}
\usepackage{placeins}
\usepackage{subcaption}
\usepackage{rotating} % sideways (landscape) floats for wide comparison tables



\usepackage{stmaryrd}
\usepackage{bookmark}
\usepackage{booktabs}
\usepackage{stackengine}
\usepackage{adjustbox}
\usepackage{footnote}
\makesavenoteenv{tabular}
\makesavenoteenv{table}


% Style de paragraphe classique LaTeX : pas d'espace entre paragraphes, alinéa
% (indentation) en début de paragraphe — identique à la référence.
\setlength{\parskip}{0pt plus 1pt}
\import{preambule}{theorems}

\input{preambule/notations}

%citation style
\let\barecite\cite
\let\cite\parencite

\import{preambule}{abbreviations}

% cleveref chargé en dernier (après hyperref, biblatex, glossaries-extra, bookmark),
% conformément à la recommandation du paquet.
\usepackage{cleveref}

% --- Correctif babel-arabic + \part/\chapter ---------------------------------
% Avec l'option « arabic » de babel, \@currentlabel ET le label interne de
% cleveref (\cref@currentlabel) sont vidés juste après \part et \chapter (mais
% pas après \section). Conséquence : \ref/\Cref{part:...} et {chap:...} rendent
% « ?? ». On enveloppe \@chapter et \@part (capturés à \AtBeginDocument, donc
% APRÈS toutes les redéfinitions de babel) pour réinscrire les deux labels.
\makeatletter
\AtBeginDocument{%
  \let\bidifix@chapter\@chapter
  \def\@chapter[#1]#2{%
    \bidifix@chapter[#1]{#2}%
    \protected@edef\@currentlabel{\csname p@chapter\endcsname\thechapter}%
    \protected@edef\cref@currentlabel{[chapter][\the\value{chapter}][]\thechapter}%
  }%
  \let\bidifix@part\@part
  \def\@part[#1]#2{%
    \bidifix@part[#1]{#2}%
    \protected@edef\@currentlabel{\csname p@part\endcsname\thepart}%
    \protected@edef\cref@currentlabel{[part][\the\value{part}][]\thepart}%
  }%
}
\makeatother